\documentclass{article}
\usepackage[utf8]{inputenc}
\usepackage{graphicx}
\usepackage{amsmath}
\usepackage{geometry}
\geometry{a4paper, margin=1in}

\begin{document}

    \chapter{Background and Literature Review}

    \section{Emotion Analysis Applications}
    Physiological signals are widely used in automated emotion detection because they reflect unconscious autonomic responses [3]. In particular, the \textbf{galvanic skin response (GSR)} – a measure of skin conductance due to sweat gland activity – is tightly linked to sympathetic nervous system arousal [4, 256]. Changes in GSR occur outside conscious control, providing an \textbf{undiluted} indicator of emotional arousal and stress [5]. This makes GSR a sensitive, objective measure in settings where self-reported emotions may be unreliable [6]. For example, \textbf{multimodal} sensing systems that combine GSR with other signals have shown high accuracy: one study found that fusing thermal imaging and photoplethysmography (PPG) signals yielded about \textbf{78.3\%} stress classification accuracy, exceeding either modality alone [7, 258]. This demonstrates the benefit of integrating multiple physiological channels for emotion recognition [8].

    The platform developed in this thesis builds on these insights [9]. It synchronises a wearable \textbf{Shimmer3 GSR+ sensor} (which samples skin conductance at 128 Hz with 16-bit resolution) with a \textbf{Topdon TC001 thermal camera} (256×192 pixels at 25 Hz) and a smartphone \textbf{RGB camera} (1920×1080 at 30 fps) [10]. All data streams are time-aligned via the Network Time Protocol, achieving synchronization within a few tens of milliseconds [11]. This hardware combination enables real-time, synchronized recording of skin conductance alongside facial thermal cues and visual features [12].

    In our experiments, we use controlled stress-inducing tasks (e.g. Stroop color-word test and the Trier Social Stress Test) where ground-truth GSR is recorded by the Shimmer sensor in parallel with the contactless data streams [13]. These paired measurements produce labeled training data for machine learning models, allowing us to predict stress levels from the remote camera-based features [14]. By capturing multiple modalities together, we ensure the data are suitable for training algorithms that infer internal emotional states from external signals [15].

    \section{Rationale for Contactless Physiological Measurement}
    Traditional GSR measurement requires direct skin contact via electrodes (typically on the fingers or palm), which has several drawbacks [17]. Contact-based GSR sensors with wires can induce motion artifacts: even slight finger movements or cable tugging can cause false spikes in the conductance signal, obscuring true stress responses [18]. Long-term use of adhesive electrodes may also cause discomfort or skin irritation [19, 260]. These issues limit the practicality of GSR monitoring outside the lab [20].

    \textbf{Contactless methods}, by contrast, aim to infer sympathetic arousal without physical contact, avoiding such artifacts [21]. Thermal imaging offers a promising contact-free alternative that sidesteps movement interference [22]. Recent work confirms the feasibility of remote monitoring: for instance, a system combining a smartphone camera for PPG with a thermal infrared camera for facial temperature achieved robust stress detection [23, 258]. They found that synchronizing both modalities gave substantially higher accuracy than either alone, highlighting the value of contactless multi-sensor fusion in stress recognition [24].

    Our platform addresses limitations seen in earlier studies by employing higher-performance sensors and tight integration [25]. The Topdon thermal camera we use provides 256×192 resolution at 25 Hz, far exceeding the 8–9 Hz frame rate of some low-cost thermal imagers used in past research [26]. Importantly, it outputs \textbf{radiometric} temperature data with ±0.1 °C accuracy per pixel, allowing precise measurement of subtle thermal changes (for example, the well-documented nose-tip cooling of only a few tenths of a degree during stress [27, 262]). Meanwhile, the 30 fps RGB video captures facial expressions and permits remote photoplethysmography [28]. Combined with the Shimmer’s high-sensitivity GSR readings, this setup is designed to eventually predict continuous GSR values from the camera-based features during controlled stressors [29]. Unlike many prior works that only perform binary stress classification (“stress” vs “no stress”), our goal is to learn a mapping to the \textbf{actual} physiological signal level, enabling real-time estimation of stress intensity [30].

    An additional motivation for exploring contactless sensing is to leverage \textbf{palm-based measures} of stress [31]. Emotional sweating is most pronounced in certain regions – notably the palms, soles, and axillae – which have a high density of eccrine sweat glands under sympathetic control [32, 256]. In fact, palmar sweating can increase in response to psychological stress even when there is no thermal cause, indicating a dedicated neural pathway for emotional sweating [33]. The palm thus offers a rich source of stress information: as sweat glands activate, skin conductance rises (the basis of GSR) and evaporative cooling can occur locally [34]. A thermal camera pointed at the hand could capture these changes as minute temperature drops or altered heat patterns on the palm [35]. Focusing on the palm may offer advantages over facial measurements because the palm’s sweat response to stress is often larger and more immediate [36]. By designing our system to accommodate imaging of the participant’s palm (for instance, by having them hold a hand in view of the thermal camera), we aim to directly capture the physiological processes that drive GSR [37].

    In summary, a contactless approach – especially one that monitors known responsive regions like the face and palms – can greatly improve comfort and reduce noise, while still tapping into the sympathetic signals critical for emotion and stress detection [38, 256].

    \section{Definitions of “Stress” (Scientific vs. Colloquial)}
    “Stress” is a multifaceted concept with differing definitions across disciplines [40]. In a classic physiological sense, Hans Selye defined stress as \textbf{“the nonspecific response of the body to any demand.”} [41] This broad, response-focused definition implies that any significant physiological arousal – whether from excitement, fear, or pain – could be considered a stress response [42]. In contrast, psychological theories (such as Lazarus and Folkman’s transactional model) emphasize the role of \textbf{cognitive appraisal}: stress is said to occur when an individual perceives a situation as threatening, challenging, or beyond their coping resources [43]. In everyday language, people often conflate stress with negative emotional states (anxiety, frustration, etc.), but scientifically it can encompass any state of heightened arousal or \textbf{allostatic load} on the body [44].

    In this thesis, we adopt a practical approach to define stress for experimental purposes [45]. We focus on \textbf{acute stressors} that reliably activate the sympathetic nervous system in a controlled way [46]. By using well-known laboratory tasks with established physiological effects (for example, the Stroop color-word interference task or the Trier Social Stress Test’s public speaking challenge), we induce transient episodes of stress in participants [47]. The \textbf{ground truth} for “stress” is thus tied to these standardized stimuli rather than to subjective self-report [48]. We primarily measure phasic skin conductance responses (SCRs) – the rapid peaks in GSR occurring \textasciitilde1–5 s after a stimulus – as indicators of the sympathetic surge [49]. Tonic skin conductance levels (the slower baseline drift) are also recorded but are less useful for identifying discrete stress events [50]. By using external triggers (stimuli) and measuring the involuntary GSR and other physiological reactions, we obtain an objective link between specific stressors and the body’s response [51]. This approach sidesteps some ambiguity in the term “stress,” since our interest is in the \textbf{physiological arousal} aspect [52]. We acknowledge, however, that our findings will be most directly applicable to acute, laboratory-induced stress [53]. Real-world or chronic stress involves additional factors (individual appraisal, context, psychological coping), which are beyond our current scope [54].

    \section{Cortisol vs. GSR as Stress Indicators}
    When evaluating stress biomarkers, one might ask why we focus on GSR (reflecting sympathetic nervous activity) instead of hormonal measures like \textbf{cortisol} [56]. Cortisol, a glucocorticoid hormone released by the hypothalamic-pituitary-adrenal (HPA) axis, is often called the “stress hormone” and can be measured noninvasively via saliva [57]. However, cortisol has significant drawbacks for real-time or immediate stress monitoring [58].

    Firstly, the cortisol response is \textbf{slower}: salivary cortisol typically peaks about \textbf{20–30 minutes} after a stressor onset [59, 264]. For example, in the Trier Social Stress Test, participants’ cortisol levels begin rising during the task but only reach maximum roughly half an hour after the stress event, long after autonomic signals like heart rate and GSR have spiked [60]. Cortisol also requires biochemical analysis (either lab assays or specialized sensors), which until recently could not be done instantaneously [61]. Even rapid assay techniques still take minutes to produce results, making cortisol unsuited for “on-line” feedback or closed-loop interventions [62]. Additionally, cortisol levels are influenced by many confounding factors – time of day (diurnal rhythm), food intake, caffeine, medication, and individual differences in HPA axis reactivity – making it a noisy measure for comparing moment-to-moment stress [63]. Large inter-individual variability means that a single cortisol reading is hard to interpret without a personalized baseline [64].

    By contrast, GSR changes are \textbf{immediate and transient}: sympathetic arousal causes the skin’s conductance to begin rising within \textbf{1–3 seconds} of a stimulus, with a clear peak often at 2–5 s and recovery within tens of seconds [65]. This near-instantaneous response allows GSR to capture rapid changes in stress or emotional arousal [66]. Typical SCR (phasic) amplitudes range from 0.1–1 µS for moderate stressors and can be easily detected by our sensor’s 128 Hz sampling (which provides fine temporal resolution) [67]. The \textbf{timing} advantage means GSR can be used in real-time monitoring or biofeedback scenarios – an important consideration for our goal of synchronized data collection for machine learning [68].

    Practically speaking, incorporating cortisol into our study would be cumbersome [69]. Collecting salivary samples during each session (e.g., via swabs or spit into tubes) and then analyzing them would disrupt the flow of experiments and yield data only after a significant delay [70]. Moreover, studies that have measured both cortisol and GSR concurrently often find only a modest correlation between them [71]. For example, simultaneous measurements during lab stressors have reported correlations on the order of \textit{r} ≈ 0.3–0.5 at best [72, 262]. This suggests that while both reflect “stress,” they may capture different aspects – cortisol reflecting the HPA axis and sustained stress, versus GSR reflecting acute sympathetic activation [73]. Indeed, cognitive stressors (mental workload) can elicit strong GSR reactions with little or no immediate cortisol change, whereas physical or social stressors engage both systems [74].

    Given these considerations, we rely on GSR as our primary stress indicator in this work [75]. Its rapid response and ease of continuous measurement make it ideal for building a real-time inference dataset, even if cortisol might be collected in other studies for complementary validation [76].

    \section{GSR Physiology and Measurement Limitations}
    To effectively use GSR data, it is important to understand both the physiology behind it and the practical limitations of measurement [78]. GSR (also known as \textbf{electrodermal activity, EDA}) is typically measured by passing a tiny current between two electrodes on the skin and measuring the conductance (or its reciprocal, resistance) [79]. The Shimmer3 GSR+ device we employ applies a constant 0.5 V excitation across two Ag/AgCl electrodes attached to the fingers (often index and middle finger) [80]. This measures skin resistance in a range from \textasciitilde10 kΩ up to \textasciitilde4.7 MΩ, which the device digitizes with 16-bit precision [81]. In terms of conductance, this corresponds to a typical range of 1–100 µS (microsiemens) for human skin, with a theoretical resolution of around 0.000076 Ω (76 µΩ) in the resistance measurement [82]. In practice, this high resolution allows the Shimmer to resolve both slow \textbf{tonic} changes (skin conductance level, SCL) and fast \textbf{phasic} bursts (skin conductance responses, SCRs) [83]. An SCR is the rapid rise and fall in conductance triggered by a sudden stimulus or moment of arousal: it usually begins \textasciitilde1–3 s after the event, peaks by \textasciitilde5 s, and recovers over \textasciitilde10–20 s [84, 262]. The ability to capture these kinetics is crucial for correlating events to physiological responses [85].

    However, GSR signals are susceptible to various artifacts and variability, which must be managed through careful experimental design and signal processing [86]. A primary issue is \textbf{motion artifact}: when a participant moves their hand or fingers, the pressure and contact at the electrodes can change, creating large spikes in the recorded conductance that are not due to sweat activity [87]. In our preliminary tests, even small movements could produce spurious deflections larger than genuine SCRs (sometimes >2 µS, compared to typical stress SCRs of 0.1–0.8 µS) [88]. We mitigate this by utilizing the Shimmer’s built-in \textbf{accelerometer} [89]. The GSR sensor package includes a 3-axis accelerometer (±16g range, also sampled at 128 Hz) [90]. We log the acceleration and flag segments of GSR data where sudden movements occur (e.g. an acceleration >0.2 g) [91]. During analysis, we can exclude or treat with caution any SCR-like peaks that coincide with motion, as these are likely artifacts [92, 262].

    Environmental factors also affect GSR [93]. Skin conductance is temperature-dependent; for instance, an ambient temperature shift of ±2 °C can drift the baseline SCL by 15–25\% (warmer skin leads to higher baseline conductance due to increased skin hydration and circulation) [93]. Similarly, humidity changes the rate of sweat evaporation and skin wetness [94]. To account for this, we record \textbf{ambient conditions} using sensors on the Shimmer and in the room: air temperature (±0.5 °C accuracy) and relative humidity (±3\%) [95]. If we observe significant swings (e.g., a sudden cooling of the room from air-conditioning), we note those periods and can later adjust or normalize the data [96, 262].

    Another practical issue is electrode gel drying [97]. The conductive gel used on electrodes can dry out over time (often after \textasciitilde45–60 minutes of recording), increasing skin-electrode impedance and causing a slow drift downward in the measured conductance [97]. In long sessions, we refresh the gel about every 30 minutes and ensure the electrodes maintain good contact in a consistent location (usually the distal phalanges of two fingers) [98, 262].

    Individual differences in GSR are notable [99]. Some people are \textbf{electrodermal “non-responders”}, estimated at about 5–10\% of the population, who show very small or no SCR peaks even under strong stimuli [99, 262]. Others have widely varying baseline SCLs (e.g. one person might sit around 2 µS, another at 20 µS) due to factors like hydration, fitness, age, or medication [100]. For example, older adults and certain clinical populations tend to have lower amplitude GSR responses [101]. These differences complicate direct comparison of raw values between subjects [102]. Our approach is to \textbf{normalize within-subject} as needed – for instance, expressing SCR amplitudes as a percentage of that individual’s prestimulus baseline level [103]. We also include a baseline rest period at the start of each recording session: a quiet 5-minute interval to establish the person’s resting SCL range [104, 262]. This provides a personal reference for what constitutes a significant SCR for that individual [105].

    In summary, to ensure high-quality GSR data for analysis, we implement several measures and processing steps:
    \begin{itemize}
        \item \textbf{Baseline recording:} Each participant sits quietly for \textasciitilde5 minutes at the start, to establish their resting conductance level and natural variability [107, 262].
        \item \textbf{Motion logging:} The Shimmer’s accelerometer is recorded alongside GSR [108]. Any trial or moment with excessive motion (beyond normal fidgeting) is either excluded or annotated as potentially artifact-laden [109, 262].
        \item \textbf{Environment logging:} We continuously log ambient temperature and humidity [110]. Sudden changes are noted so we can later model or remove any baseline shifts caused by the environment [111, 262].
        \item \textbf{Signal filtering:} A light filter is applied in software – for example, a 3-sample median filter – to remove isolated electrical spikes or dropouts, while preserving the shape of genuine SCR waves [112, 262].
        \item \textbf{Electrode maintenance:} For long experiment sessions, electrode gel is replenished every 30 minutes and electrode placement is kept consistent to prevent drift [113, 262].
    \end{itemize}

    Finally, our stress induction protocol is designed to respect the physiological \textbf{refractory periods} of GSR [114]. SCRs will habituate or even overlap if stress stimuli are presented in rapid succession [115]. To avoid this, we space stimuli (for example, Stroop test trials or surprise events) about 20–30 seconds apart, ensuring that one SCR has time to peak and mostly recover before the next stimulus arrives [116, 262]. This timing preserves the full amplitude of each response and yields cleaner data where individual SCR peaks can be uniquely associated with specific events [117]. Taken together, these precautions and techniques help maximize the integrity of the GSR signals, providing a reliable ground truth for our contactless sensing experiments [118].

    \section{Thermal Cues of Stress in Humans}
    One of the main contactless approaches for stress sensing involves \textbf{thermal infrared imaging} of the skin [120]. Stress-induced thermal patterns on the face and body have been documented in prior research [121, 262]. The physiological link is twofold: (1) changes in blood flow due to sympathetic vasoconstriction, and (2) heat loss due to sweat evaporation [122]. Under acute stress, the SNS redirects blood flow away from the skin (a fight-or-flight response to preserve core circulation), which can cause peripheral areas to cool [123]. For example, studies using high-resolution thermal cameras have observed that cognitive stress tasks produce a detectable cooling of the nose tip by about \textbf{0.3–0.5 °C} [124, 262]. The nose is an extremity with dense blood vessels and little muscle, so a reduction in blood perfusion shows up clearly as a temperature drop [125]. In one recent report, the magnitude of nose-tip cooling (on the order of a few tenths of a degree) correlated with the person’s self-reported stress levels (correlation \textit{r} ≈ 0.6–0.7) [126, 262]. This makes nasal temperature a useful proxy for stress-induced sympathetic activity [127].

    Our equipment is well-suited to capture these subtle thermal cues [128]. The Topdon TC001 thermal camera provides a full radiometric image at 25 Hz with sensitivity around 0.1 °C [129]. This is sufficient to detect the rapid vasoconstriction that begins within seconds of stress onset in regions like the nose [130]. Importantly, because the camera gives actual temperature readings for each pixel (instead of just a false-color image), we can quantify changes precisely (e.g., nose skin cooling from 33.0 °C to 32.5 °C) [131]. We focus on specific \textbf{regions of interest (ROIs)} on the face that are known to reflect stress-related autonomic changes [132]:
    \begin{itemize}
        \item \textbf{Nose tip (and alae):} As noted, vasoconstriction in the nose causes a characteristic cooling during stress. We track a small ROI at the nasal tip, where the effect tends to be strongest [133, 134, 262].
        \item \textbf{Forehead/periorbital region:} Some studies report sympathetic \textbf{warming} in the forehead or areas around the eyes under mental stress (possibly due to increased metabolism or muscle tension). We include a forehead ROI to monitor this, as a comparison to the nose cooling [135, 136, 262].
        \item \textbf{Temperature gradients:} Rather than absolute temperature at one spot, the \textbf{difference} between regions can be informative. Stress often produces a “cold center, warm periphery” pattern on the face [137, 138]. We compute gradients like nose-to-forehead or nose-to-cheek temperature difference to see this pattern (the nose cooling while perhaps the cheeks stay warmer) [139, 262].
        \item \textbf{Breathing signals (nostrils):} Thermal imaging can also capture respiration. As the person exhales warm air from the nostrils, the area just in front of the nose periodically warms [140, 141]. Under stress, breathing rate and depth might change (often becoming faster and shallower) [142]. By monitoring the cyclic thermal fluctuations at the nostrils, we can extract respiratory rate and relative depth; slower, deeper breathing versus rapid panting can be an indicator of stress or relaxation [143, 144, 262].
    \end{itemize}
    To reliably detect these thermal signals, we control the recording environment [145]. The lab is kept at a stable room temperature (\textasciitilde22 ± 1 °C) so that moderate vasomotor changes are not overshadowed by external temperature shifts [146]. We also use lighting that emits minimal infrared (LED lights) to avoid accidentally heating the subject’s face with lamps [147]. Participants sit about 0.8–1.2 m from the thermal camera, such that key facial features (like the nose) span \textasciitilde10 pixels or more – enough for a good thermal reading [148]. Before each task, we record a 2-minute thermal baseline while the subject is at rest, creating an individualized reference thermal map [149, 262]. During the stress task, we then look at \textbf{changes relative to baseline} for each ROI (e.g., how much the nose temp drops from that person’s resting value) [150]. We also cross-validate these thermal events with GSR peaks; for instance, a pronounced nose-tip cooling occurring simultaneously with a large SCR would strongly indicate a true stress response [151].

    It’s worth noting that thermal imaging captures some signals that are \textbf{involuntary and impossible to consciously mask} [152]. Unlike facial expressions (which can be faked or suppressed) or even heart rate to an extent (which can be altered by breathing techniques), thermal signatures like vasoconstriction or sweat-based cooling cannot be directly controlled by the subject [153]. This gives thermal sensing a robustness in detecting genuine stress [154]. Prior research in psychophysiology has found that thermal infrared imaging (sometimes termed \textbf{functional infrared imaging, fITI}) can be as sensitive as GSR for detecting arousal [155, 256]. In some cases, thermal imaging even differentiated subtle affective states, providing additional information beyond a binary “stress vs. no stress” signal [156, 256].

    While our primary thermal focus is on the face (due to its rich innervation and convenience of imaging), we recognize that other body regions exhibit strong thermal responses too [157]. Palmar areas, for instance, might show temperature changes due to sweating (evaporative cooling) and altered blood flow when a person is stressed [158]. Some studies have directly visualized sweat gland activation on the fingers using high-speed thermal cameras [159, 266]. In this work, we did not image the palms during the main protocol, but in future experiments the \textbf{palm could be targeted} as well [160]. A thermal view of the palm might reveal sweat-induced cooling spots or changes in perfusion that correlate with the GSR signal recorded on that same hand [161]. Since the palm is a principal site of emotional sweating, we anticipate that thermal features from the hand (e.g., the appearance of cooler, damp areas) could provide an even more direct index of sympathetic activation [162]. Integrating such palm thermal data with facial thermal and GSR would further strengthen a multi-region approach to contactless stress sensing [163].

    \section{RGB vs. Thermal Imaging for Machine Learning}
    Given the complementary nature of visual and thermal cues, a central hypothesis of this research is that \textbf{fusing RGB video with thermal infrared data} will improve stress detection performance over either modality alone [165]. Prior studies provide a basis for this hypothesis [166]. Under controlled conditions, thermal imaging of the face alone has achieved high stress classification accuracies (in some reports, exceeding 80\% for distinguishing stress vs. calm) by leveraging features like temperature changes and breathing signals [166, 258]. Visible-light video (RGB) of the face, on the other hand, can capture signals such as facial expressions, heart rate via remote PPG, and head/body movements [167]. However, approaches using RGB video alone typically reach more modest accuracies (often in the 65–75\% range for binary stress classification in lab studies) [168, 258]. This gap is because facial expressions can be consciously controlled or may not reliably indicate internal stress, and RGB-based physiological measures are sensitive to lighting and motion [169]. Thus, each modality has limitations: \textbf{thermal} excels at sensing involuntary physiological changes (but thermal cameras have lower resolution and may miss behavioral context), while \textbf{RGB} provides rich behavioral information (but can be fooled by intentional masking or poor lighting) [170]. We expect that combining these modalities can yield a more robust system – essentially, each modality can cover the other’s “blind spots.” [171].

    Concretely, we plan to train machine learning models on three sets of features: (a) thermal-only features, (b) RGB-only features, and (c) a fusion of thermal + RGB features [172]. The target output for these models will be the continuous GSR value or a derivative metric, treating it as an indicator of stress level [173]. By comparing model performance, we can test if \textbf{multimodal fusion} truly offers an advantage in predicting stress (we anticipate it will) [174]. This approach is inspired by findings like those of Cho et al. (2019), who showed that combining thermal nose-tip data with smartphone PPG significantly improved acute stress inference accuracy (78\% vs \textasciitilde69\% with PPG alone) [175, 176, 258]. In our case, thermal data will contribute features related to autonomic changes (temperature drops, breathing rate changes), while RGB data will contribute features related to expressions and cardiovascular activity (remote pulse) [177].

    On the RGB side, we extract several features using state-of-the-art computer vision techniques [178]. We employ face tracking (e.g., MediaPipe or similar) to obtain facial landmarks and track expression changes (frowns, brow furrows, etc.) [179]. This gives quantitative measures of facial muscle movements that might indicate stress (for instance, furrowing of the brow or tensing of facial muscles around the eyes) [180]. We also implement a \textbf{remote PPG (rPPG)} algorithm on the RGB video: typically, the green channel of the video is analyzed over skin regions (like the cheeks or forehead) to detect the subtle periodic color changes caused by blood pulsing [181]. From this we can derive the heart rate and even heart rate variability (HRV) of the participant [182]. HRV is known to decrease under stress (due to increased sympathetic and reduced parasympathetic activity) [183]. By capturing rPPG, we effectively get a proxy for cardiovascular stress responses in a contactless way [184]. That said, RGB signals have known limitations: they can be corrupted by changes in ambient light, shadows on the face, or head motions that break the tracking [185]. Moreover, facial expressions are partially under voluntary control – a participant might suppress visible signs of stress (a neutral expression) even while their physiology is in turmoil [186].

    Thermal imaging provides a complementary set of features that are largely involuntary [187]. As described in Section 2.6, we obtain features like the nose-tip temperature trajectory, the forehead-nose gradient, and breathing rhythm [188]. These are directly tied to autonomic responses that cannot be faked or hidden [189]. The thermal camera is also indifferent to lighting conditions – it works in darkness or bright light since it senses emitted heat [190]. This means it can operate in scenarios where an RGB camera might fail (e.g., low-light conditions or where color information is ambiguous) [191]. On the other hand, thermal images are lower resolution and lack the detailed textures or color cues of RGB; they cannot capture, say, a subtle smile or a change in skin tone (blushing) that an RGB camera might [192, 193].

    By fusing the two, we hope to leverage the \textbf{strengths of both}: the thermal modality catching the physiological “ground truth” signals of stress and the RGB modality catching contextual and complementary cues [194]. All data streams in our system are tightly synchronized, which is crucial for effective sensor fusion [195]. Each thermal frame and RGB frame is timestamped at capture, and our synchronization (via NTP and shared triggers) keeps them aligned to within \textasciitilde20 ms [196]. We validate this alignment by cross-checking known synchronous events – for example, a deep breath should appear as a simultaneous thermal change (at the nose) and a subtle chest movement (in the RGB video) [197]. Indeed, we included a calibration event (like a quick facial temperature change using a cold stimulus visible on camera) to ensure both cameras record the event at the same time [198]. With alignment confirmed, we can fuse features frame-by-frame or on a common time base for the model [199].

    In summary, our \textbf{working hypothesis} is that a combined model using both RGB and thermal features will outperform single-modality models for estimating stress (as indexed by GSR) [200]. Prior evidence from literature supports this idea: for instance, the combination of remote PPG signals with nasal thermal signals was shown to infer stress better than either alone [201, 258]. By rigorously testing this hypothesis on our synchronized dataset, we aim to quantify the improvement from sensor fusion [202]. A successful outcome would be a significant increase in prediction accuracy or correlation with true GSR when using the fused feature set, confirming that the two modalities provide complementary information about the user’s affective state [203].

    \section{Sensor Hardware Selection Rationale}
    A key contribution of this project is the careful selection and integration of sensors to build a \textbf{multimodal stress-sensing platform} [205]. We chose specific hardware components based on their capabilities, compatibility, and the demands of high-quality data collection [206]:
    \begin{itemize}
        \item \textbf{Shimmer3 GSR+ (Wearable GSR/PPG Sensor):} This device was selected as the \textbf{ground-truth physiological sensor} for several reasons [207]. It is a research-grade wearable that records GSR at 128 Hz with 16-bit ADC resolution, covering a wide resistance range suitable for EDA (10 kΩ–4.7 MΩ) [208]. This high resolution ensures even minute conductance changes are captured [209]. The Shimmer3 also includes a PPG module (for pulse) and a 3-axis accelerometer, which we leverage for artifact rejection (as discussed in Section 2.5) [210]. The unit communicates via Bluetooth and comes with an open API, simplifying real-time data streaming into our custom software [211]. Shimmer devices are well-validated in psychophysiology research, lending confidence that the GSR data is accurate and comparable to published studies [212]. We considered alternatives (such as the Empatica E4 wristband or BIOPAC’s wired systems), but the Shimmer’s combination of \textbf{high sampling rate, data fidelity, and open integration} best met our needs [213]. It also has a battery life of \textasciitilde10–12 hours, allowing long recording sessions [214]. We use standard Ag/AgCl gel electrodes from the manufacturer to ensure consistency with calibration settings [215]. Overall, the Shimmer3 provides a reliable GSR reference and additional signals (like PPG, motion) in one compact unit [216].

        \item \textbf{Topdon TC001 Thermal Camera (Smartphone-Mounted):} For the thermal imaging component, we chose the Topdon TC-series infrared camera, which is a lightweight, portable module that attaches to an Android smartphone via USB-C [217]. It offers a thermal resolution of 256×192 at 25 Hz and – critically – outputs \textbf{radiometric data} for each pixel [218]. Many inexpensive thermal cameras (FLIR One, SEEK Thermal) have lower frame rates (9 Hz) or do not provide raw temperature readings, only images [219]. The Topdon camera’s specs (frame rate, sensitivity \textasciitilde0.1 °C) meet the requirements identified in literature for detecting small stress-related temperature changes [220]. At a cost of only a few hundred USD, it is far more affordable than laboratory-grade thermal cameras (which can cost tens of thousands), making our solution scalable and reproducible [221]. Topdon provides an SDK (InfiSense) that allows us to adjust settings (such as emissivity correction, ranges) and to stream the raw thermal frames into our Android app [222]. We selected this device over similar smartphone thermal cameras because of its \textbf{full SDK support and raw data access} – many low-end cameras lock users into proprietary apps or provide only processed images [223]. The Topdon also has an onboard battery, so it does not heavily drain the phone, and it adheres to the UVC standard (USB Video Class), which improves compatibility [224]. In summary, this camera hits a sweet spot of performance vs. cost, enabling high-quality thermal data collection in a mobile setup [225, 262].

        \item \textbf{Samsung Galaxy S22 Smartphone (RGB Camera):} We use a modern Android smartphone (Samsung S22, released 2022) as the RGB video sensor and as the host for the thermal camera [226]. The S22 was chosen for its high-quality built-in camera and strong hardware support for imaging [227]. It can record 1080p video at 30 fps (or higher if needed) with a large camera sensor that performs well in low light [228]. Crucially, the S22 supports \textbf{Camera2 API Level 3}, meaning we can capture \textbf{RAW} images and have manual control over exposure and focus [229]. By using the phone’s \textbf{Expert RAW} mode or Camera2 API, we avoid the usual smartphone processing (compression, auto white-balance) that could distort subtle color changes needed for rPPG [230]. Capturing uncompressed, 10-bit or 12-bit RAW frames preserves fine details (such as slight skin color fluctuations) that a typical 8-bit webcam or phone video might miss [231, 268]. The smartphone’s computing power also allows us to run real-time processing (like face tracking) on the device [232]. Using an off-the-shelf smartphone has practical benefits: it makes the platform more \textbf{portable and accessible}, and demonstrates that consumer devices can be repurposed for sophisticated physiological sensing [233]. In fact, earlier work by Cho et al. (2019) showed that a smartphone camera plus a thermal addon can successfully detect stress [234, 258], supporting our decision to use a phone-based system [234]. The Galaxy S22 specifically gives us a robust development environment (running Android 13 with support for our custom app) and can easily be mounted or carried as needed [235].
    \end{itemize}

    \subsection*{Synchronization and Integration}
    All the sensors above are orchestrated to act as a single system [236]. We run a \textbf{PC-based controller} (a Python program) that communicates with the smartphone and the Shimmer [237]. The Shimmer GSR device streams its data via Bluetooth to the PC [238]. Simultaneously, the Android phone (with the Topdon attached) runs our app that captures thermal frames and RGB frames [239]. These frames are timestamped on the phone and sent to the PC over a Wi-Fi network socket in real time [240]. To ensure a common time base, we operate an NTP server on the PC and have the phone periodically sync its clock to it, achieving time consistency on the order of ±10–20 ms [241]. We also log events or markers (like task start/stop signals) on both devices [242]. If needed, we use a visual/audible marker (for example, a flash of an LED in view of both cameras) to post-synchronize and verify timing [243]. The end result is that all data – GSR, thermal video, RGB video, and any other streams (accelerometer, etc.) – share a unified timeline [244, 262]. This level of synchronization is critical for machine learning later on, as it allows aligning the input features with the correct output labels (GSR) without ambiguity [245].

    In summary, the chosen sensor suite (a high-grade GSR unit, a mid-resolution thermal imager, and a modern smartphone camera) provides \textbf{high-quality, synchronized data} for our multimodal stress sensing application [246]. Each component was selected with guidance from prior research and our specific engineering requirements [247]. For instance, the importance of raw signal fidelity and timing led us to the S22 phone and Shimmer device, while the demonstrated utility of mobile thermal sensing [248, 258, 262] guided our choice of the Topdon camera [248]. By combining these into an integrated platform, we aim to produce a dataset and system that meet research-grade standards for accuracy and temporal precision [249]. This allows us to explore the core research questions under realistic conditions and paves the way for future real-time stress inference systems built on the foundation of this hardware [250]. Each hardware choice is furthermore backed by related work [251]. The suitability of smartphones for PPG and thermal monitoring has been demonstrated by multiple groups (e.g., measuring stress with phone camera HR and thermal nose data [251, 258]). Shimmer sensors have been used as a \textbf{gold standard} in many affective computing studies for EDA, thanks to their reliability and open data access [252]. The Topdon camera’s specifications align with recommendations from thermal imaging studies (which often note that >0.1 °C sensitivity and ≥25 Hz frame rate are needed to capture quick autonomic changes) [253]. Our integration approach draws inspiration from open-source toolkits that emphasize synchronized multi-sensor recording [254, 256]. By leveraging and extending these insights, the assembled platform is well-justified for the research at hand and positioned to significantly advance contactless emotion sensing capabilities [255].

    \begin{thebibliography}{99}
        \itemsep0em
        \bibitem{1} Placeholder for: Thermal infrared imaging in psychophysiology: Potentialities and limits - PMC. Cited for insights on sympathetic nervous system arousal, thermal imaging as sensitive as GSR, and toolkits for multi-sensor recording.
        \bibitem{2} Placeholder for: JMIR Mental Health - Instant Stress: Detection of Perceived Mental Stress Through Smartphone Photoplethysmography and Thermal Imaging. Cited for stress classification accuracy using fused thermal and PPG signals, and the feasibility of smartphone-based stress detection.
        \bibitem{3} Placeholder for: A LaTeX file (1.tex). Cited regarding the potential for skin irritation from long-term use of adhesive electrodes.
        \bibitem{4} Placeholder for: Chapter 2\_ Background and Literature Review (1).docx. Cited multiple times for details on nose-tip cooling during stress, correlation of cooling with self-reported stress, SCR kinetics, motion artifact mitigation, environmental effects on GSR, electrode maintenance, electrodermal non-responders, baseline recording protocols, stimulus spacing, and facial ROIs for thermal imaging.
        \bibitem{5} Placeholder for: Thermal infrared imaging in psychophysiology: Potentialities and limits - PMC. Cited for the high density of eccrine sweat glands in palmar regions.
        \bibitem{6} Placeholder for: A meta-analysis of cortisol reactivity to the Trier Social Stress Test. Cited for the typical 20-30 minute peak time of salivary cortisol after a stressor.
        \bibitem{7} Placeholder for: Chapter 2\_ Background and Literature Review (1).docx. Cited for modest correlation values between cortisol and GSR during lab stressors.
        \bibitem{8} Placeholder for: Chapter 2\_ Background and Literature Review (1).docx. Cited for modest correlation values between cortisol and GSR during lab stressors.
        \bibitem{9} Placeholder for: Chapter 2\_ Background and Literature Review (1).docx. Cited for the timing of skin conductance responses (SCRs) after a stimulus.
        \bibitem{10} Placeholder for: Chapter 2\_ Background and Literature Review (1).docx. Cited for the use of an accelerometer to identify and exclude motion artifacts in GSR data and for adjusting data based on ambient conditions.
        \bibitem{11} Placeholder for: Chapter 2\_ Background and Literature Review (1).docx. Cited for noting environmental changes for data normalization and applying median filters to GSR signals.
        \bibitem{12} Placeholder for: Chapter 2\_ Background and Literature Review (1).docx. Cited for the practice of replenishing electrode gel during long sessions.
        \bibitem{13} Placeholder for: Chapter 2\_ Background and Literature Review (1).docx. Cited for ensuring good electrode contact in a consistent location.
        \bibitem{14} Placeholder for: Chapter 2\_ Background and Literature Review (1).docx. Cited for the concept of electrodermal "non-responders".
        \bibitem{15} Placeholder for: Chapter 2\_ Background and Literature Review (1).docx. Cited for the use of a baseline rest period to establish a personal SCL range.
        \bibitem{16} Placeholder for: Chapter 2\_ Background and Literature Review (1).docx. Cited for the exclusion or annotation of trials with excessive motion.
        \bibitem{17} Placeholder for: Chapter 2\_ Background and Literature Review (1).docx. Cited for spacing stimuli to respect the refractory period of SCRs.
        \bibitem{18} Placeholder for: Chapter 2\_ Background and Literature Review (1).docx. Cited for tracking ROIs at the nasal tip and forehead.
        \bibitem{19} Placeholder for: Chapter 2\_ Background and Literature Review (1).docx. Cited for computing temperature gradients between facial regions.
        \bibitem{20} Placeholder for: Chapter 2\_ Background and Literature Review (1).docx. Cited for extracting respiratory rate from thermal fluctuations at the nostrils.
        \bibitem{21} Placeholder for: Chapter 2\_ Background and Literature Review (1).docx. Cited for recording a thermal baseline before tasks.
        \bibitem{22} Placeholder for: Thermal infrared imaging in psychophysiology: Potentialities and limits - PMC. Cited for thermal imaging differentiating subtle affective states.
        \bibitem{23} Placeholder for: Thermal infrared imaging in psychophysiology: Potentialities and limits - PMC. Cited for thermal imaging providing more than a binary stress signal.
        \bibitem{24} Placeholder for: Infrared thermal imaging for assessing human perspiration and evaluating antiperspirant product efficacy | Scientific Reports. Cited for the visualization of sweat gland activation using high-speed thermal cameras.
        \bibitem{25} Placeholder for: Chapter 2\_ Background and Literature Review (1).docx. Cited in the context of the Topdon camera's performance versus cost.
        \bibitem{26} Placeholder for: Chapter 2\_ Background and Literature Review (1).docx. Cited in the context of the Topdon camera's suitability for a mobile setup.
        \bibitem{27} Placeholder for: Sensor Device Selection Rationale.docx. Cited regarding the benefits of capturing RAW image frames for preserving fine details.
        \bibitem{28} Placeholder for: Sensor Device Selection Rationale.docx. Cited regarding the benefits of capturing uncompressed frames over typical processed video.
        \bibitem{29} Placeholder for: Chapter 2\_ Background and Literature Review (1).docx. Cited for the detail that all data streams share a unified timeline.
        \bibitem{30} Placeholder for: Chapter 2\_ Background and Literature Review (1).docx. Cited for the detail that all data streams share a unified timeline.
        \bibitem{31} Placeholder for: JMIR Mental Health - Instant Stress: Detection of Perceived Mental Stress Through Smartphone Photoplethysmography and Thermal Imaging. Cited for the demonstrated utility of mobile thermal sensing.
    \end{thebibliography}

\end{document}